Entscheiden Sie, ob die folgenden Funktionen eine elementare Stammfunktion
haben und finden Sie sie, falls sie existiert.
\begin{teilaufgaben}
\item
$\displaystyle \int \frac{dx}{x\log(x)^2}$
\item
$\displaystyle \int \frac{dx}{x^2\log(x)}$
\end{teilaufgaben}

\begin{loesung}
Beide Teilaufgaben verlangen nach einer Stammfunktion in $\mathbb{Q}(x,\vartheta)$ mit $\vartheta=\log x$.
Der Integrand ist ein rationaler Teil mit $a(\vartheta)=1$.
\begin{teilaufgaben}
\item Für $b(\vartheta)=x\vartheta^2$ wird die Methode von Rothstein-Trager
angewendet.
Dazu müssen nichttriviale gemeinsame Teiler von
\[
a(\vartheta)-c\cdot Db(\vartheta)
=
1-c\cdot (\vartheta + 1)
\qquad\text{und}\qquad
x\vartheta^2
\]
gefunden werden.
Nur für $c_1=1$ bleibt vom linken Term nur $\vartheta$, welches auch
gleichzeitig der grösste gemeinsame Teiler $v_1(\vartheta)=\vartheta$ ist.
Es folgt
\[
\int \frac{dx}{x\log(x)^2}
=
\log v_1(\vartheta)
=
\log(\log(x)).
\]
\item
In diesem Fall ist $b(\vartheta)=x^2\vartheta$ mit $Db(\vartheta) = 2x\vartheta + x$.
Gesucht sind nichttriviale gemeinsame Teiler von
\[
a(\vartheta)-c\cdot Db(\vartheta)
=
1-c\cdot (2\vartheta+1)x
=
(1-cx) -2c\vartheta
\qquad\text{und}\qquad
x^2\vartheta.
\]
Ein nichttrivialer gemeinsamer Teiler ist nur möglich, wenn $1-cx=0$ ist
oder $c=1/x$.
Da $c$ nicht konstant ist, ist das Integral nicht elementar.
\qedhere
\end{teilaufgaben}
\end{loesung}
